\uuid{nTkS}
\titre{Solution générale/raccord}
\auteur{hanani ; mzari}
\organisation{exo7}
\licence{CC BY-NC-SA 4.0}
\datecreate{2020-01-20}
\niveau{L1}
\module{Analyse}
\chapitre{Equation différentielle}
\sousChapitre{Résolution d'équation différentielle du premier ordre}
\difficulte{3}
\theme{Equations du premier ordre, solution générale, raccord}
\exo7qcmid{lille-2.equations-differentielles.014}

\contenu{
\begin{question}
\qclassification{225.01}
\qtags{motcle=solution générale/raccord}
On considère l'équation différentielle $(E)$ : $t^2y'=y$. Parmi les
affirmations suivantes, cocher celles qui sont vraies :
\begin{answers}
\good{$(E)$ est une équation linéaire homogène du premier ordre.}
\bad{La fonction $\displaystyle y=\mathrm{e}^{1/t}$ est une solution de $(E)$ sur $]0,+\infty[$.}
\good{La solution générale de $(E)$ sur $]0,+\infty[$ est $\displaystyle y=k\mathrm{e}^{-1/t}$, $k\in \Rr$.}
\bad{La fonction nulle est l'unique solution de $(E)$ sur $\Rr$.}
\end{answers}
\begin{explanations}
Une primitive de $-1/t^2$ est $1/t$, donc la solution générale de $(E)$, sur $]0,+\infty[$ (ou sur $]-\infty ,0[$) est
$$y=k\mathrm{e}^{-1/t},\quad k\in \Rr.$$
Si $y$ est une solution de $(E)$ sur $\Rr$, alors $y(0)=0$ et
$$y(t)=\left\{\begin{array}{ll}k_1\mathrm{e}^{-1/t}&\mbox{si }t>0\\ k_2\mathrm{e}^{-1/t}&\mbox{si }t<0.\end{array}\right.$$
La continuité en $0$ implique que $k_2=0$. Ainsi
$$y(t)=k_1\mathrm{e}^{-1/t}\mbox{ si }t>0\mbox{ et }y(t)=0\mbox{ si }t\leq 0.\leqno{(\star)}$$
De plus, une telle fonction est dérivable en $0$. Réciproquement, toute fonction définie par $(\star)$ est une solution de $(E)$ sur $\Rr$.
\end{explanations}
\end{question}
}
